\section{Lean formalization architecture}

The artifact's \pathref{lean/FORMALIZATION_PLAN.md} gives a file-by-file
implementation schedule.  This section records the mathematical dependency
structure and the intended theorem signatures.

\subsection{Module graph}

\begin{description}[leftmargin=3.4cm,style=nextline]
\item[\leanname{FQNF/KernelCore}.]
Logical proof bundles, simultaneous-cut compilation, reflected checker
specialization, QF loop extraction, independent-certificate stability, and
boundary/gap adapters.  This file contains no placeholder.

\item[\leanname{FQNF/BoolReflection}.]
Boolean arithmetic, reflected equality, propositional closure, bounded products
and sums, and compilation of bounded formulas.

\item[\leanname{FQNF/Digits}.]
Power places, base digits, array encoders, mixed-radix prefix products, pack/get
operations, code bounds, and decoding theorems.

\item[\leanname{FQNF/CertificateSearch}.]
Conversion from one-valued verifiers to the zero-valued convention of
\leanname{Search.mu}; explicit-witness and existential selection theorems.

\item[\leanname{FQNF/PathCode}.]
Path-state decoding, bounded path checkers, whole-path selection, and the
stepwise seriality adapter.

\item[\leanname{FQNF/Query}.]
The high-level package structure and the main compiler theorem.

\item[\leanname{FQNF/Names}.]
Rational intervals, cross-precision stability, demand transformers, and
apartness-guarded exact operations.
\end{description}

\subsection{Structural core}

The supplied core defines
\[
 \mathsf{Derives}(\Gamma,\Delta)
  :\equiv \forall A\in\Delta\;\Gamma\objectvdash A.
\]
Its central theorem has the form
\begin{lstlisting}[language=Lean]
theorem compileLocal
    (hDelta : Derives Gamma Delta)
    (hlocal : Delta ++ Gamma |- [PRAAx] phi) :
    Gamma |- [PRAAx] phi
\end{lstlisting}
and is a direct application of \leanname{Prf.trans}.  A reflected checker
stores fixed syntax and closed soundness/completeness derivations.  Specializing
those derivations at arbitrary object terms uses the existing
\leanname{IPRA.inst} theorem.

This core should be elaborated first because it tests imports and notation
without depending on any new arithmetic.

\subsection{Digit arithmetic proof order}

The most substantial generic formalization is the digit theorem.  The safe
order is:

\begin{enumerate}
\item prove positivity of powers under \(b\ge2\);
\item define prefix sums and prove the partial-code bound;
\item prove the lower/current/upper decomposition of an encoded array;
\item apply \leanname{div\_mod\_unique} at \(b^i\);
\item apply it again at \(b\);
\item generalize powers to mixed-radix prefix products;
\item prove substitution and congruence lemmas for every new term constructor.
\end{enumerate}

The decomposition is the point at which informal formalizations often rely on
ring reasoning over semantic naturals.  The object proof should instead use the
repository's \leanname{ipra\_ring}, bounded-sum index lemmas, and explicit
\(\Prf\) congruence.

\subsection{Selection adapter}

The repository search theorem already proves the hard induction.  The new
adapter needs only:

\begin{enumerate}
\item Booleanity of \(V\);
\item \(V=1\to 1\mono V=0\);
\item \(1\mono V=0\to V=1\) under Booleanity;
\item specialization of \leanname{Search.search\_found};
\item an existential-elimination proof whose conclusion contains the fixed
selected term rather than the temporary witness.
\end{enumerate}

This should be implemented before any generic viability table.  Whole-path
search then gives a correct fallback for every bounded finite relation.

\subsection{Moving checked glue into the derivation library}

Two relevant source files currently live under proof packets:
\pathref{ObjectIteration.lean} and \pathref{ObjectObservation.lean}.  Their
reported status is stronger than most packet blueprints: they construct actual
\(\Prf\,\PRAAx\) terms, and the latter records a successful compilation audit.
For long-term reuse, their generic theorems should be moved or rederived under
\code{IPRADerivations/Core} or \code{IPRADerivations/FQNF}, with imports adjusted
and axiom checks registered.

\subsection{Axiom and placeholder audit}

A completed integration should add every public theorem to
\pathref{Derivations/Lean/IPRADerivations/Checks/Axioms.lean}.  The acceptance
criteria are:

\begin{itemize}
\item no \code{sorryAx}, \code{admit}, or new logical axiom;
\item conclusions are literal \code{Prf PRAAx} statements;
\item host-side decidability or normalization is used only in the metatheory;
\item every object computation has defining equations or a reflected checker;
\item every package README distinguishes checked declarations from executable
reference tests.
\end{itemize}

\subsection{Recommended first end-to-end theorem}

The finite path example from \cref{thm:whole-path} is the best integration
benchmark.  It exercises digits, bounded Boolean folds, search, existential
elimination, and extraction, but requires no difficult domain mathematics.
After that, formalize the exact rational cubic bisection example from
\cref{sec:examples}; it adds a packed trace and a nontrivial arithmetic
invariant.
